\documentclass[11pt]{article}

\usepackage{amsmath, amssymb, amsthm}
\usepackage{mathtools}
\usepackage[colorlinks=true]{hyperref}
\usepackage[capitalize]{cleveref}
\usepackage{pdflscape}
\usepackage{tikz}
\usetikzlibrary{cd}

\theoremstyle{definition}
\newtheorem{definition}{Definition}[section]
\newtheorem{theorem}[definition]{Theorem}
\newtheorem{lemma}[definition]{Lemma}
\newtheorem{remark}[definition]{Remark}

\newcommand{\id}{\mathrm{id}}
\newcommand{\outp}[1]{#1^{+}}
\newcommand{\inpt}[1]{#1^{-}}
\newcommand{\ev}{\mathrm{ev}}
\newcommand{\coev}{\mathrm{coev}}
\newcommand{\To}[1]{\xrightarrow{#1}}
\newcommand{\Fun}[1]{\textsf{#1}}
\newcommand{\cat}[1]{\mathcal{#1}}
\newcommand{\ihom}[2]{[#1,#2]}

\begin{document}

\section{Setup}

Let $(\cat{C},I,\otimes,\ihom{-}{-})$ be a symmetric monoidal closed category with evaluation $\ev\colon\ihom{A}{B}\otimes A\to B$. Let $(\Fun{T},\eta,\mu,\sigma)$ be a strong monoidal monad on $\cat{C}$ with strength $\sigma_{M,N}\colon M\otimes\Fun{T}N\to\Fun{T}(M\otimes N)$, and let $(z,\alpha\colon\Fun{T}z\to z)$ be a $\Fun{T}$-algebra. Write $H_c\coloneqq\ihom{\inpt{c}}{z}$.

A coKleisli lens $f\colon\binom{\inpt{c}}{\outp{c}}\to\binom{\inpt{d}}{\outp{d}}$ has:
\begin{itemize}
\item $\outp{f}\colon\outp{c}\to\outp{d}$ (a comonoid homomorphism), and
\item $\inpt{f}\colon\outp{c}\otimes\inpt{d}\to\Fun{T}\inpt{c}$.
\end{itemize}

\begin{definition}[$\psi_f$]\label{def.psi}
Define $\psi_f\colon\outp{c}\otimes H_c\to H_d$ to be the composite
\[
\begin{tikzcd}[column sep=40pt, row sep=20pt, ampersand replacement=\&]
\outp{c}\otimes H_c
  \ar[r, "{\coev}"]
  \ar[dd, dashed, "\psi_f"']
\& \ihom{\inpt{d}}{\outp{c}\inpt{d} H_c}
  \ar[d, "{\ihom{\inpt{d}}{\inpt{f}\otimes\id}}"]
\\
\& \ihom{\inpt{d}}{\Fun{T}\inpt{c}\; H_c}
  \ar[d, "{\ihom{\inpt{d}}{\sigma}}"]
\\
H_d
\& \ihom{\inpt{d}}{\Fun{T}(\inpt{c}\otimes H_c)}
  \ar[d, "{\ihom{\inpt{d}}{\Fun{T}\ev}}"]
\\
\& \ihom{\inpt{d}}{\Fun{T}z}
  \ar[ul, "{\ihom{\inpt{d}}{\alpha}}"]
\end{tikzcd}
\]
That is: $\coev$, then apply $\ihom{\inpt{d}}{-}$ to each of $\inpt{f}\otimes\id$, $\sigma$, $\Fun{T}\ev$, $\alpha$.
\end{definition}

\begin{definition}[$\Theta_z$]
$\Theta_z(f)\colon\outp{c}\otimes H_c\to\outp{d}\otimes H_d$ is the composite
\[
\outp{c}\otimes H_c
\To{\delta\otimes\id}
\outp{c}\otimes\outp{c}\otimes H_c
\To{\outp{f}\otimes\psi_f}
\outp{d}\otimes H_d.
\]
\end{definition}

\section{The lemma}

\begin{lemma}\label{lem.psi_compose}
For composable coKleisli lenses $f\colon\binom{\inpt{c}}{\outp{c}}\to\binom{\inpt{d}}{\outp{d}}$ and $g\colon\binom{\inpt{d}}{\outp{d}}\to\binom{\inpt{e}}{\outp{e}}$:
\[
\psi_{g\circ f}=\psi_g\circ(\outp{f}\otimes\psi_f)\circ(\delta_{\outp{c}}\otimes\id_{H_c}).
\]
\end{lemma}

The right side is:
\[
\outp{c}\otimes H_c
\To{\delta\otimes\id}
\outp{c}\otimes\outp{c}\otimes H_c
\To{\outp{f}\otimes\psi_f}
\outp{d}\otimes H_d
\To{\psi_g}
H_e.
\]

\section{Step-by-step proof in the monad case}

Both sides of \cref{lem.psi_compose} are maps $\outp{c}\otimes H_c\to H_e=\ihom{\inpt{e}}{z}$. We use juxtaposition for $\otimes$.

\subsection{The left side: $\psi_{g\circ f}$ spelled out}

Following \cref{def.psi} applied to $g\circ f$, the map $\psi_{g\circ f}\colon\outp{c} H_c\to H_e$ is:
\[
\begin{tikzcd}[column sep=35pt, row sep=20pt, ampersand replacement=\&]
\outp{c} H_c
  \ar[r, "{\coev}"]
\& \ihom{\inpt{e}}{\outp{c}\inpt{e} H_c}
  \ar[d, "{\ihom{\inpt{e}}{\delta}}"]
\\
\& \ihom{\inpt{e}}{\outp{c}\outp{c}\inpt{e} H_c}
  \ar[d, "{\ihom{\inpt{e}}{\outp{f}}}"]
\\
\& \ihom{\inpt{e}}{\outp{c}\outp{d}\inpt{e} H_c}
  \ar[d, "{\ihom{\inpt{e}}{\inpt{g}}}"]
\\
\& \ihom{\inpt{e}}{\outp{c}\Fun{T}\inpt{d}\; H_c}
  \ar[d, "{\ihom{\inpt{e}}{\sigma}}"]
\\
\& \ihom{\inpt{e}}{\Fun{T}(\outp{c}\inpt{d})\; H_c}
  \ar[d, "{\ihom{\inpt{e}}{\Fun{T}\inpt{f}}}"]
\\
\& \ihom{\inpt{e}}{\Fun{T}\Fun{T}\inpt{c}\; H_c}
  \ar[d, "{\ihom{\inpt{e}}{\mu}}"]
\\
\& \ihom{\inpt{e}}{\Fun{T}\inpt{c}\; H_c}
  \ar[d, "{\ihom{\inpt{e}}{\sigma}}"]
\\
\& \ihom{\inpt{e}}{\Fun{T}(\inpt{c} H_c)}
  \ar[d, "{\ihom{\inpt{e}}{\Fun{T}\ev}}"]
\\
\& \ihom{\inpt{e}}{\Fun{T}z}
  \ar[d, "{\ihom{\inpt{e}}{\alpha}}"]
\\
\& H_e
\end{tikzcd}
\]
Every arrow after $\coev$ applies $\ihom{\inpt{e}}{-}$ to a single morphism. The first six (from $\delta$ to $\mu$) are $\ihom{\inpt{e}}{\inpt{(g\circ f)}\otimes\id_{H_c}}$; the last three are shared with every $\psi$.

\subsection{The right side: $\psi_g\circ(\outp{f}\otimes\psi_f)\circ(\delta\otimes\id)$ spelled out}

This is three maps composed. The first is $\delta\otimes\id\colon\outp{c} H_c\to\outp{c}\outp{c} H_c$. The second applies $\outp{f}$ to one copy of $\outp{c}$ and $\psi_f$ (from \cref{def.psi}) to the other copy with $H_c$:
\[
\begin{tikzcd}[column sep=35pt, row sep=20pt, ampersand replacement=\&]
\outp{c}\outp{c} H_c
  \ar[r, "{\outp{f}\otimes\coev}"]
\& \outp{d}\;\ihom{\inpt{d}}{\outp{c}\inpt{d} H_c}
  \ar[d, "{\ihom{\inpt{d}}{\inpt{f}\otimes\id}}"]
\\
\& \outp{d}\;\ihom{\inpt{d}}{\Fun{T}\inpt{c}\; H_c}
  \ar[d, "{\ihom{\inpt{d}}{\sigma}}"]
\\
\& \outp{d}\;\ihom{\inpt{d}}{\Fun{T}(\inpt{c} H_c)}
  \ar[d, "{\ihom{\inpt{d}}{\Fun{T}\ev}}"]
\\
\& \outp{d}\;\ihom{\inpt{d}}{\Fun{T}z}
  \ar[d, "{\ihom{\inpt{d}}{\alpha}}"]
\\
\& \outp{d}\; H_d
\end{tikzcd}
\]
The third applies $\psi_g$ (same pattern as \cref{def.psi} with $g$ in place of $f$):
\[
\begin{tikzcd}[column sep=35pt, row sep=20pt, ampersand replacement=\&]
\outp{d} H_d
  \ar[r, "{\coev}"]
\& \ihom{\inpt{e}}{\outp{d}\inpt{e} H_d}
  \ar[d, "{\ihom{\inpt{e}}{\inpt{g}\otimes\id}}"]
\\
\& \ihom{\inpt{e}}{\Fun{T}\inpt{d}\; H_d}
  \ar[d, "{\ihom{\inpt{e}}{\sigma}}"]
\\
\& \ihom{\inpt{e}}{\Fun{T}(\inpt{d} H_d)}
  \ar[d, "{\ihom{\inpt{e}}{\Fun{T}\ev}}"]
\\
\& \ihom{\inpt{e}}{\Fun{T}z}
  \ar[d, "{\ihom{\inpt{e}}{\alpha}}"]
\\
\& H_e
\end{tikzcd}
\]
The right side of \cref{lem.psi_compose} composes: $\delta\otimes\id$, then the first diagram, then the second.

\subsection{The proof diagram}

Both sides of \cref{lem.psi_compose} begin with $\coev_e$; by naturality of $\coev_e$, the comparison reduces to showing two composites $\outp c\otimes\inpt e\otimes H_c\to z$ agree inside $\ihom{\inpt e}{-}$. Both composites further share the opening $\delta$, $\outp f$ (using cocommutativity of $\delta_{\outp c}$), so the diagram below starts from $\outp c\otimes\outp d\otimes\inpt e\otimes H_c$ inside $\ihom{\inpt e}{-}$, with juxtaposition for~$\otimes$ and writing $c$ for $\inpt c$, etc.\ as in the preceding sections. The top-right path is the right side of \cref{lem.psi_compose}; the left-bottom path is the left side.

\begin{landscape}
\thispagestyle{empty}
\[
\resizebox{\linewidth}{!}{%
\begin{tikzcd}[column sep=0.4em, row sep=1.2em, every label/.append style={font=\tiny}, font=\tiny, ampersand replacement=\&]
%% Row 0: before g^-, RHS psi_f path
\outp c\outp d e H_c
  \ar[r, "{\coev}"]
  \ar[d, "{\inpt g}"']
\& \outp d e{[d,\outp c d H_c]}
  \ar[r, "{\inpt f}"]
  \ar[d, "{\inpt g}"']
\& \outp d e{[d,\Fun T(c) H_c]}
  \ar[r, "{\sigma}"]
  \ar[d, "{\inpt g}"']
\& \outp d e{[d,\Fun T(c H_c)]}
  \ar[r, "{\ev}"]
  \ar[d, "{\inpt g}"']
\& \outp d e{[d,\Fun Tz]}
  \ar[r, "{\alpha}"]
  \ar[d, "{\inpt g}"']
\& \outp d e H_d
  \ar[d, "{\inpt g}"]
\\
%% Row 1: after g^-
\outp c\Fun T(d) H_c
  \ar[d, "{\sigma}"']
  \ar[r, "{\coev}"]
\& \Fun T(d){[d,\outp c d H_c]}
  \ar[r, "{\inpt f}"]
  \ar[d, "{\sigma}"']
\& \Fun T(d){[d,\Fun T(c) H_c]}
  \ar[r, "{\sigma}"]
  \ar[d, "{\sigma}"']
\& \Fun T(d){[d,\Fun T(c H_c)]}
  \ar[r, "{\ev}"]
  \ar[d, "{\sigma}"']
\& \Fun T(d){[d,\Fun Tz]}
  \ar[r, "{\alpha}"]
  \ar[d, "{\sigma}"']
\& \Fun T(d) H_d
  \ar[d, "{\sigma}"]
\\
%% Row 2: after sigma (inside T)
\Fun T(\outp c d) H_c
  \ar[dd, "{\inpt f}"']
  \ar[dr, "{\sigma}"]
\& \Fun T(d{[d,\outp c d H_c]})
  \ar[r, "{\inpt f}"]
  \ar[d, "{\ev}"]
\& \Fun T(d{[d,\Fun T(c) H_c]})
  \ar[r, "{\sigma}"]
  \ar[dd, "{\ev}"']
\& \Fun T(d{[d,\Fun T(c H_c)]})
  \ar[r, "{\ev}"]
  \ar[dd, "{\ev}"']
\& \Fun T(d{[d,\Fun Tz]})
  \ar[r, "{\alpha}"]
  \ar[dd, "{\ev}"']
\& \Fun T(d H_d)
  \ar[dd, "{\ev}"]
\\
%% Row 3: T(ev_d) landing for column 1 only
\& \Fun T(\outp c d H_c)
  \ar[dr, "{\inpt f}"]
\&\&\&\&
\\
%% Row 4: after T(f^-) / T(ev_d)
\Fun T\Fun T(c) H_c
  \ar[d, "{\mu}"']
  \ar[rr, "{\sigma}"]
\&\& \Fun T(\Fun T(c) H_c)
  \ar[r, "{\sigma}"]
\& \Fun T(\Fun T(c H_c))
  \ar[r, "{\ev}"]
  \ar[d, "{\mu}"']
\& \Fun T(\Fun Tz)
  \ar[r, "{\alpha}"]
  \ar[d, "{\mu}"']
\& \Fun Tz
  \ar[d, "{\alpha}"]
\\
%% Row 5: after mu, final
\Fun T(c) H_c
  \ar[rrr, "{\sigma}"]
\&\&\& \Fun T(c H_c)
  \ar[r, "{\ev}"]
\& \Fun Tz
  \ar[r, "{\alpha}"]
\& H_e
\end{tikzcd}
}
\]
\end{landscape}

Each rectangular tile commutes by one of: functoriality of $\otimes$ (rows~0--1); naturality of $\sigma$ (rows~1--2, columns~1--5); naturality of $\ev$ inside~$\Fun{T}$ (rows~2--4, columns~1--5); naturality of $\mu$ (rows~4--5, columns~3--4); or the $\Fun{T}$-algebra axiom $\alpha\circ\mu=\alpha\circ\Fun{T}\alpha$ (rows~4--5, columns~4--5).

Three non-rectangular regions remain. The pentagon at rows~1--3, columns~0--1 (the $\sigma$ diagonal from $(2{,}0)$ to $(3{,}1)$) commutes because $\Fun{T}\ev\circ\sigma\circ(\coev\otimes\id)=\sigma$ by the zig-zag identity of the $({-}\otimes d)\dashv\ihom{d}{{-}}$ adjunction inside~$\Fun{T}$, together with naturality of $\sigma$ at $\coev$. The parallelogram with vertices $(2{,}0)$, $(3{,}1)$, $(4{,}0)$, $(4{,}2)$ commutes by naturality of $\sigma$ at $\inpt{f}$. The large tile at rows~4--5, columns~0--3 commutes by $\mu$-coherence of the strength: $\sigma\circ(\mu\otimes\id)=\mu\circ\Fun{T}\sigma\circ\sigma$.

\end{document}
